\section{Unified approach for potential-field transformations}

Given a particular equivalent layer $p_{\varepsilon}$ (Definition \ref{def:continuous-eqlayer}) 
defined as a square-integrable function by equation \ref{eq:p-Fourier}, in terms of specific $\alpha_{d}$ and $\alpha_{a}$, 
it is theoretically possible to compute any potential field transformation (equations \ref{eq:transformed-field-space} and 
\ref{eq:transformed-field-Fourier}).
It means that, contrarily to the common sense, we do not need to use a planar layer of dipoles to compute the RTP anomaly from 
the approximated total-field anomaly, for example. In fact, even a layer of monopoles can be used to perform this potential-field transformation. 
This can be verified by first noting that applying the Fourier-convolution theorem to \ref{eq:transformed-field-continuous-layer} 
on a planar surface $\mathcal{S}^{\ddagger}$ (Figure \ref{fig:Figure1}) gives
\begin{equation}
	\mathcal{F} [ f_{\varepsilon}(\mathbf{r}^{\ddagger}) ] =
	\mathcal{F} \left[ p_{\varepsilon}(\mathbf{r}'') \right]
	\mathcal{F} [ t(\mathbf{r}^{\ddagger}, \mathbf{r}'') ] \: , 
	\quad \mathbf{r}^{\ddagger} \in \mathcal{S}^{\ddagger} \: ,
	\quad \mathbf{r}'' \in \mathcal{S}'' \: ,
	\label{eq:transformed-field-continuous-layer-Fourier}
\end{equation}
where
\begin{equation}
	\mathcal{F} [ t(\mathbf{r}^{\ddagger}, \mathbf{r}'') ] = \alpha_{t} \, \mathcal{F} [ \Psi(\mathbf{r}^{\ddagger}, \mathbf{r}'') ]
	\label{eq:kernel-t-Fourier}
\end{equation}
is defined in terms of a differential operator $\alpha_{t}$ acting on $\Psi$ (equation \ref{eq:Psi}) in the Fourier domain.
Then we set 
\begin{equation}
	\alpha_{t} = \left( \tfrac{\alpha_{a}}{\alpha_{d}} \right) \alpha_{f} \: ,
	\label{eq:alpha-t-generic}
\end{equation}
where $\alpha_{f}$ is the differential operator associated with the desired potential field $f$ (equation \ref{eq:transformed-field-Fourier}).
In this case, we get from equations \ref{eq:Psi-Fourier}, \ref{eq:p-Fourier}  and \ref{eq:kernel-t-Fourier} 
that the transformed potential field $f_{\varepsilon}$ (equation \ref{eq:transformed-field-continuous-layer-Fourier}) 
is given by:
\begin{equation}
	\mathcal{F} [ f_{\varepsilon}(\mathbf{r}^{\ddagger}) ] =
	\alpha_{f} \, k \, \mathcal{F} [ U_{\varepsilon}(\mathbf{r}'') ] \: 
	\mathcal{F} [ \Psi(\mathbf{r}^{\ddagger}, \mathbf{r}'') ] \: , 
	\quad \mathbf{r}^{\ddagger} \in \mathcal{S}^{\ddagger} \: ,
	\quad \mathbf{r}'' \in \mathcal{S}'' \: .
	\label{eq:transformed-field-continuous-layer-Fourier2}
\end{equation}
Finally, we use \ref{eq:Psi-Fourier}, \ref{eq:transformed-field-Fourier} and \ref{eq:potential-U-epsilon-Fourier-layer} to show that
\begin{equation}
	\mathcal{F} [ f_{\varepsilon}(\mathbf{r}^{\ddagger}) ] = 
	\begin{cases}
		\mathcal{F} [ f(\mathbf{r}^{\ddagger}) ] &
		, \quad 0 < k \le k_{\varepsilon} \\
		0 & , \quad k > k_{\varepsilon} 
	\end{cases} \: ,
	\label{eq:transformed-field-continuous-layer-1-Fourier}
\end{equation}
if the auxiliary potential $U_{\varepsilon}$ is defined by equation \ref{eq:potential-U-epsilon-Fourier-dagger-1},
or
\begin{equation}
	\mathcal{F} [ f_{\varepsilon}(\mathbf{r}^{\ddagger}) ] = 
	\tfrac{1}{1 + \mu_{\varepsilon} \, k^{2} \, e^{\, k \, (z''_{0} - z^{\ddagger}_{0})}} \; 
	\mathcal{F} [ f(\mathbf{r}^{\ddagger}) ] \: ,
	\label{eq:transformed-field-continuous-layer-2-Fourier}
\end{equation}
if the auxiliary potential $U_{\varepsilon}$ is defined by equation \ref{eq:potential-U-epsilon-Fourier-dagger-2}.
In both cases, we can see that the transformed field $f_{\varepsilon}$ yielded by the equivalent layer 
is a smooth approximation of that produced by the true sources on the planar surface $\mathcal{S}^{\ddagger}$
(Figure \ref{fig:Figure1}). Note that, to obtain a correct potential-field transformation leading to an $f_{\varepsilon}$ 
that approximates $f$, we need to correctly set the operator $\alpha_{t}$ (equation \ref{eq:alpha-t-generic}), 
regardless of the equivalent layer type defined by $\alpha_{a}$.


In the context of the equivalent-layer technique, a potential-field transformation is generally defined
in the space domain, by following two steps. The first consists in obtaining a particular layer $p_{\varepsilon}$ 
(Definition \ref{def:continuous-eqlayer}), which is subsequently used to compute a transformed field $f_{\varepsilon}$ via 
\ref{eq:transformed-field-continuous-layer}.
It can be verified, however, that this two-step approach is equivalent to the single-step approach based on the 
Fourier transform of the observed potential field (equation \ref{eq:field-transformation-Fourier}). 
We find by using \ref{eq:observed-data-Fourier} and \ref{eq:transformed-field-continuous-layer-Fourier2} that
\begin{equation}
	\mathcal{F} [ f_{\varepsilon}(\mathbf{r}^{\ddagger}) ] =
	\left( \tfrac{\alpha_{f}}{\alpha_{d}} \right) \, \mathcal{F} [ d_{\varepsilon}(\mathbf{r}^{\ddagger}) ]
	 \: .
	\label{eq:transformed-field-continuous-layer-single-step}
\end{equation}
Then, if we consider that the predicted field $d_{\varepsilon}$ (equation \ref{eq:predicted-data-continuous-layer-epsilon-Fourier})
is defined in terms of the auxiliary potential $U_{\varepsilon}$ given by equation \ref{eq:potential-U-epsilon-Fourier-dagger-1},
we get
\begin{equation}
	\mathcal{F} [ f_{\varepsilon}(\mathbf{r}^{\ddagger}) ] =
	\left( \tfrac{\alpha_{f}}{\alpha_{d}} \right) \, 
	\begin{cases}
		\mathcal{F} \left[ d^{o}(\mathbf{r}^{\dagger}) \right] 
		e^{\, k \, (z^{\ddagger}_{0} - z^{\dagger}_{0})} & , \quad 0 < k \le k_{\varepsilon} \\
		0 & , \quad k > k_{\varepsilon} 
	\end{cases} \: .
	\label{eq:transformed-field-continuous-layer-1-single-step}
\end{equation}
Alternatively, if we consider that the predicted field $d_{\varepsilon}$ (equation \ref{eq:predicted-data-continuous-layer-epsilon-Fourier})
is defined in terms of the auxiliary potential $U_{\varepsilon}$ given by equation \ref{eq:potential-U-epsilon-Fourier-dagger-2},
we get
\begin{equation}
	\mathcal{F} [ f_{\varepsilon}(\mathbf{r}^{\ddagger}) ] =
	\left( \tfrac{\alpha_{f}}{\alpha_{d}} \right) \, 
	\tfrac{1}{1 + \mu_{\varepsilon} \, k^{2} \, e^{\, k \, (z''_{0} - z^{\dagger}_{0})}} \;
	\mathcal{F} \left[ d^{o}(\mathbf{r}^{\dagger}) \right] \: .
	\label{eq:transformed-field-continuous-layer-2-single-step}
\end{equation}
Note that, for the particular case in which $z^{\ddagger}_{0}$ coincides with $z^{\dagger}_{0}$ and 
the truncating wavenumber $k_{\varepsilon}$ goes to infinity, equation 
\ref{eq:transformed-field-continuous-layer-1-single-step} reduces to \ref{eq:field-transformation-Fourier}.
Similarly, if $z^{\ddagger}_{0}$ coincides with $z^{\dagger}_{0}$ and $\mu_{\varepsilon}$ goes to zero,
equation \ref{eq:transformed-field-continuous-layer-2-single-step} also reduces to \ref{eq:field-transformation-Fourier}.
For the first time, we provide a mathematical proof establishing the link between
potential-field transformations performed using the equivalent-layer technique in 
the space domain and those computed via Fourier transformation of the observed data.
This connection sheds light on why potential-field transformations based on the
equivalent layers can exhibit the same numerical instabilities commonly associated
with Fourier-based transformations of potential-field data. This
relationship is widely acknowledged and is typically justified, without any proof, by  
the notion that physical phenomena are independent of the domain 
-- whether spatial or Fourier -- used for their computation.